\begin{enumerate}

\item El metal ideal determina la componente tangencial absoluta.

\[\fitmath{\mathrm{ctid}=\mathrm{uu}+\mathrm{cm}\cdot \tan\left(\mathrm{lam}\right)}\]

\[\fitmath{\mathrm{ctid}=(448)+(139)\cdot \tan\left((-0.5061454831)\right)\approx 370.951\,\mathrm{m/s}}\]

\item Se aplica el cociente de remolinos definido.

\[\fitmath{\mathrm{ct}=\mathrm{slip}\cdot \left(\mathrm{uu}+\mathrm{cm}\cdot \tan\left(\mathrm{lam}\right)\right)}\]

\[\fitmath{\mathrm{ct}=(0.89)\cdot \left((448)+(139)\cdot \tan\left((-0.5061454831)\right)\right)\approx 330.1464\,\mathrm{m/s}}\]

\item Sin prerrotación, Euler queda U2 Ctheta2.

\[\fitmath{\mathrm{wc}=\mathrm{uu}\cdot \mathrm{slip}\cdot \left(\mathrm{uu}+\mathrm{cm}\cdot \tan\left(\mathrm{lam}\right)\right)}\]

\[\fitmath{\mathrm{wc}=(448)\cdot (0.89)\cdot \left((448)+(139)\cdot \tan\left((-0.5061454831)\right)\right)\approx 147905.6\,\mathrm{J/kg}}\]

\end{enumerate}
